\subsection{GENERAL OFFICERS}

General Officer counters are provided inthe SHENANDOAH Game for all of the most important officers who served there during the times covered by the various scenarios. Not all generals are included, just the most prominent ones. These officers had a great influence on the campaigns and battles that took place in the Shenandoah Valley, their skill and ability in handling their units often being the difference between victory and defeat.

Their ability in various important functions has been evaluated, and is simulated in the Game through the Leadership Value number found on each General Officer counter (the larger the number, the more effective the commander). A quick glance at the Leadership Values of these various counters reveals that, although the Confederacy committed fewer troops to the Valley than did the Union, they did some of their finest generals to command them. The difference in leadership led to many humiliating reverses for the generally superior Union forces in the Valley.

\begin{enumerate}
  \item General Officer counters "commadn" the units they are in the same hex with. Their effect will be only on these units.

  \item General Officer Counters have no Combat Factors. If they are in an enemy Zone of Control, and not Stacked with any friendly Combat Factors, they are automatically eliminated.

  \item The Ashby Confederate General Officer counter has no effect on units other than cavalry and Horse Artillery. The Sheridan counter can affect all types of units. All other General Officer Counters can affect only infantry and artillery units.

  \item In a case where several General Officer counters are in the same hex, only one can be used at a time; excess counters are placed in the "Reserve", and play no part.

  \item In a case wehre units from several hexes are attacking the same enemy hex, only one General Officer counter (of the player's choice) in the attack can be used for the Combat Modifier. In no case will the Leadership Value numbers of General Officer counters be added together.

  \item \textbf{COMBAT:}

  General Officer counters can modify the Basic Combat Table of their units. Compare the Leadership Value numbers of the opposing commanders, subtracting the smaller number from the larger one (if there is no opposing General Officer counter, this equals zero). This number can be added to the Basic Combat Table number of the force with the higher rated commander.

  \item \textbf{RALLYING:}

  A player with a General Officer counter in the hex with units which have just suffered a "Disordered" result in Combat may attempt to negate this "Disorder" by rolling a die and consulting the TROOP RALLYING TABLE (note that this can result in success, failure, or failure plus the elimination of the General's counter). If successful, the "Disorder" result can be ignored.

  \item \textbf{FORCED MARCHES:}

  The Phasing Player may attempt to "Force March" units Stacked with the General Officer Counter at any time during his Phase. Units which "Force March" are allowed to use two more Movement Allowance Factors than normal (i.e. an Infantry unit is increased to eight Movement Allowance Factors, a Cavalry unit to twelve Movement Allowance Factors, etc.)

  The Combat penalties (see the Forch March column of the TABLE MODIFIERS chart) apply only while the two extra Movement Allowance Factors are being used at the end of the Phase, not while the regular ones are being used earlier in the Phase. The player attempting the "Force March" may do so if a die is rolled and the result equals the General's Leadership Value number, or less. Otherwise, it cannot be done. Note that it is possible for a General Officer Counter to be with several different Stacks during a Turn, and each could roll to Force March while there.
\end{enumerate}